%!TEX root = ../../main.tex
% PPT绘图说明：左侧画一条包含成功、失败、重试和空闲区间的原始episode时间轴。
% 中间画prepare operators，将同一episode拆成完整episode、segment和transition三种粒度。
% 右侧画进入learner的training units，并为不同unit标注training role和sampling weight。
% 底部用一条细线表示lineage，使每个unit可以追溯到原始episode及其变换步骤。
\begin{figure*}[t]
\centering
\fbox{\rule{0pt}{5.2cm}\rule{0.94\textwidth}{0pt}}
\caption{Planned illustration of how one raw episode can yield training units at different granularities and with different training roles. The final artwork should retain the provenance link from each unit to the corresponding interval in the raw episode.}
\label{fig:data-unit-transformation}
\end{figure*}
